% 19-limitations.tex: current limitations and what external validation is required.
\section{Current limitations and the external validation still required}
\label{sec:limitations}

The repository's readiness ledger opens with one line, and this paper repeats it as its own bound:
at \code{1.0.0-rc.7} Polaris is not production-ready for real identity data. Every engineering gap
the ledger enumerated is closed and pinned by a check; what remains is not buildable in the
repository. The list below separates what has not been done from what cannot be done by the author,
and it is longer than Version 2's in one direction, because the repository spent the interval
finding limits rather than removing them.

\subsection{Decisions only a deploying organisation can make}

Nine decisions must be recorded for a named deployment before real data could be held: a legal
basis, impact assessment and regulator approval; the signing-key custody driver and who holds the
key; the high-availability topology and its hosts, and whether a second region is run; encryption at
rest and its key custodian; the off-site backup target and schedule; the alerting backend, the
on-call rotation and who receives the duress page; the right-to-erasure policy; the retention
schedule per class, with the counsel who says the number satisfies the statute; and an independent
penetration test with a threat-model sign-off. Polaris records each decision and its justification;
it does not know the law. \sExt

\subsection{What is not built}

\begin{itemize}
  \item The physical token is a specification, an emulator and published vectors; no card is
    manufactured and no silicon is certified, and the profile states which of its properties an
    emulator cannot establish.
  \item No hardware security module has been used; the PKCS\#11 driver is proven against a software
    token, the KMS driver against a stand-in, and the lattice library is not FIPS-validated.
  \item A second region is a standby cluster with a non-zero recovery point; status distribution is
    a set of cache rules, not a deployment; and the national workload has been simulated at a scale
    factor, never run.
  \item A hash-based signer is registered and not wired. The hop from the pooler to the database,
    the certificate signatures and today's operator authenticators remain classical.
  \item A verifier cannot learn that an algorithm has been deprecated, so retiring one is a software
    release to every relying party; rollback and the mixed window of a partial migration are
    unmeasured.
  \item Blinded and one-time presentations are not built, so a plain credential shows a verifier a
    stable token value, an issuer signature and a holder key, and an SD-JWT VC presentation shows nine
    issuer-signed values that are the same everywhere. The per-relying-party handle bounds what a
    verifier \emph{stores}; only the zero-knowledge path bounds what it is \emph{shown}, and that
    bound is the size of the epoch, which nothing floors above one and no verdict reports.
  \item Neither the scoped nullifier nor the pairwise handle hides a holder from the \emph{issuer},
    which derives every epoch leaf to build the tree. Issuer-blind derivation is not claimed.
  \item Polaris as an issuer has no attribute credential and no anonymous predicate; as a verifier
    it accepts a disclosed boolean the issuer had to anticipate.
  \item The duress mechanism resists the casual coercer and neither the informed coercer beside the
    holder nor lawful or institutional access, against which its indelible record is net-negative.
    No holder-side panic signal exists and none is proposed.
  \item Document authenticity at enrolment is taken on the operator's word; the assurance level is
    derived from what evidence was presented, not from whether it was genuine.
  \item Of the identifier columns that grow at the enrolment rate, two remain 32-bit,
    \code{Individual.individual\_id} and \code{IdentityToken.token\_id}: about twenty-nine years of
    headroom at the stated surge of 200,000 a day, past the capacity model's horizon; the interim credential during a recovery cool-down is not built; a sunset verdict
    rests on operator attestations the system cannot verify.
\end{itemize}

\subsection{What has been done by someone else, and what has not}

\begin{itemize}
  \item One unmodified wallet written by someone else has presented a credential and been accepted,
    on one path, in one format, with one signature algorithm, and it refused Polaris once first.
    No second wallet, no wallet on the native protocol, and no outside relying party.
  \item The OpenID Foundation's hosted suite ran the verifier profile: eleven modules, zero failures,
    seven negative modules passed by the service's own scoring, four positive modules in review
    awaiting a human the Foundation supplies at submission. Not submitted, not certified.
  \item A published third-party corpus agrees with both signature witnesses on the ML-DSA-65
    primitive. It tests the primitive, not the protocol.
  \item Four packages are on public registries, each verified by installing from the registry into a
    clean environment. Nobody outside the repository has reported installing one, and a registry
    download count is not a person.
  \item No external penetration test, cryptographic audit, red team, bug bounty or certification
    has occurred. The scope document and the review packet exist; no engagement has been
    commissioned. Almost every guarantee in this paper is checked by machinery written by the same
    author as the guarantee.
  \item No operator other than the author has run any of it, which is the one condition that
    separates the release candidate from \code{1.0.0}. No pilot with consenting users has run; the
    data has always been notional.
  \item No external team has implemented the native protocol from the specification and conformance
    suite alone; the suite is the contract and the integration has not happened.
  \item No institutionally independent witness watches any log; the witnesses in the record are
    processes in a drill. Two federated instances are two instances of one code base run by one
    author.
  \item The circuit over the proof library has had no external review, and the concrete bit
    security of the shipped proof configuration has not been independently derived.
  \item The repository's thesis, that a stranger can read it cold and correctly infer what it
    needs, was never tested by a stranger before its deadline and is recorded as inconclusive, with
    the strong claim permanently retired; the stranger's path of Section~\ref{sec:use} is the
    narrower test that replaced it, and nobody outside has yet reported walking it.
\end{itemize}

\subsection{What the numbers are}

Every performance number in this paper is single-node on one laptop, or a simulation at a stated
scale factor, or a projection that the repository labels as one under a named assumption. No
production cluster has been measured, and the one measurement on a real failover topology is a
median at two requests per second with the tail withheld because the sample cannot carry it. The
national planning targets are targets; the arithmetic that clears them is the easy half of the
question, and the half that nearly failed was an integer column.

\subsection{Documentation that lags the code}

Researching this version found places where the repository's own prose is behind its code or its
history; they are recorded in the author's notes that accompany this version rather than fixed
silently, as Version 2's were. Where this paper and any repository document disagree, the paper
followed the code at \code{1.0.0-rc.7} and says so in the note.
